\chapter{真空能量、暴脹與原初漲落}

\section{真空能量與德西特空間}

在平坦空間中，弗里德曼方程為
\begin{equation}
\left(\frac{\dot A}{A}\right)^2=\frac{8\pi G}{3}\rho.
\end{equation}
對物質而言，\(\rho\)按\(A^{-3}\)的比例下降；對輻射而言，\(\rho\)按\(A^{-4}\)的比例下降。真空能量則不會稀釋：
\begin{equation}
\rho=\rho_0=\text{constant}.
\end{equation}
如果宇宙中只有這種真空能量，方程右側便是常數。將這個常數寫成\(H^2\)，
\begin{equation}
H^2=\frac{8\pi G}{3}\rho_0,
\qquad
H=\sqrt{\frac{8\pi G}{3}\rho_0}.
\end{equation}
哈勃常數便是真正的常數：
\begin{equation}
\frac{\dot A}{A}=H.
\end{equation}
因此
\begin{equation}
\dot A=HA,
\end{equation}
解呈指數形式：
\begin{equation}
A(t)=C e^{Ht}.
\end{equation}
彼此遙遠且不受束縛的星系之間，距離會按指數增長。這就是德西特空間。

在這些座標下，時空幾何為
\begin{equation}
d\tau^2=dt^2-e^{2Ht}\left(dx^2+dy^2+dz^2\right).
\end{equation}
此處的真空能量為正。在平坦空間中，它必須為正，因為弗里德曼方程左側非負。負的真空能量只有在納入曲率項時才可能出現，而這裡討論的並非那種情況。

愛因斯坦引入宇宙常數，原本是為了另一個目的：平衡普通物質，讓宇宙維持靜態。

\section{當今加速膨脹的宇宙}

德西特空間也是當今宇宙晚期演化的趨勢。較早時，標度因子遵循物質主導時的關係
\begin{equation}
A(t)\propto t^{2/3}.
\end{equation}
幾十億年前，曲線開始向上彎，增長速度也逐漸加快。由於仍有其他形式的能量存在，膨脹還不是完美的指數形式，但趨勢正朝此方向發展。

這裡所說的加速是指
\begin{equation}
\ddot A>0.
\end{equation}
這是宇宙晚期膨脹的一項觀測特徵。

現在來看視界。哈勃定律表示
\begin{equation}
v=HD.
\end{equation}
在退行速度等於光速的距離處，
\begin{equation}
D=\frac{c}{H}.
\end{equation}
等價地，也可以改寫為\(HD=c\)。\editorialnote{這是將口頭說明的關係式\(D=c/H\)以代數方式重新表述。圖~\ref{fig:lecture07-horizon}右側黑板最上方那一行只有部分可辨。}

\begin{figure}[ht]
\centering
\includegraphics[width=0.62\textwidth]{lecture_07_figure_02.png}
\caption{德西特膨脹中的宇宙視界關係\(D=c/H\)。}
\label{fig:lecture07-horizon}
\lectureframe{00:30:37}
\end{figure}

這個距離約為一百億至一百五十億光年；若將宇宙年齡用\(c\)換算成距離，數值也相近。

\begin{classroomqa}
\lecturetimestamp{00:14:10}
\audiencequestion{視界尺度與宇宙年齡幾乎相等，難道不正是切線恰好通過原點時會出現的情況嗎？}
\lecturerresponse{若是\(A(t)\propto t\)，那就完全成立。但一般而言並非如此。在德西特空間中，\(H\)維持不變，宇宙年齡卻持續增加，因此目前兩者幾乎相等並非數學上的恆等關係。}
\end{classroomqa}

\section{真空能量的意義}

這裡以普通量子力學的語言描述真空能量。諧振子的能量概略可寫成
\begin{equation}
E\sim p^2+x^2.
\end{equation}
在古典力學中，基態就是平衡點：
\begin{equation}
x=0,\qquad p=0,\qquad E=0.
\end{equation}
量子力學不允許\(x\)和\(p\)同時精確為零。不確定性原理使系統必然有漲落，因此基態能量為正。諧振子的基態能量為
\begin{equation}
E_0=\frac{1}{2}\hbar\omega.
\end{equation}

接著將同樣的想法應用於場。宇宙中充滿諧振子：電磁場的振盪，以及其他每一種場的振盪。每種場都有零點能。電磁場的能量密度概略可寫成
\begin{equation}
u_{\mathrm{EM}}\sim E^2+B^2.
\end{equation}
即使沒有真實光子，仍存在無法消除的零點振盪。這是真空能量的一種解釋。此處只需要知道，這種能量確實存在，而且不會隨宇宙膨脹而稀釋。

\begin{classroomqa}
\lecturetimestamp{00:16:15}
\audiencequestion{真空能量究竟是什麼？如果移除物質和輻射，為什麼仍有能量留下？}
\lecturerresponse{因為真空是量子場的基態，而量子基態的能量不一定為零。諧振子的基態能量為正，而場可以視為諧振子的集合。因此，即使沒有真實粒子，仍有零點能存在。}
\end{classroomqa}

這不是三度背景輻射。微波背景是普通輻射，其能量密度會下降，因為光子數密度降低，而且每個光子都會紅移。真空能量則是移除真實光子後仍留下的零點能。

\section{\(H^{-1}\)的意義}

現在回到倒數哈勃常數與宇宙年齡之間的關係。如果
\begin{equation}
A(t)\propto t,
\end{equation}
那麼
\begin{equation}
H=\frac{\dot A}{A}=\frac{1}{t}.
\end{equation}
更一般地，若
\begin{equation}
A(t)\propto t^p,
\end{equation}
那麼
\begin{equation}
H=\frac{p}{t}.
\end{equation}
對物質主導的情況\(p=2/3\)而言，
\begin{equation}
H=\frac{2}{3t}.
\end{equation}
因此，在冪律宇宙模型中，\(H^{-1}\)與宇宙年齡同量級，但不會恰好等於宇宙年齡。

德西特空間則不同。在那裡
\begin{equation}
A(t)=C e^{\gamma t},
\end{equation}
因此
\begin{equation}
H=\frac{\dot A}{A}=\gamma,
\end{equation}
是個真正的常數。此處的\(H^{-1}\)不是宇宙年齡，而是指數增長的時間常數：標度因子增為原來\(e\)倍所需的時間。口頭講解中一再說``加倍''，隨即又更正；精確的說法是增為\(e\)倍。如今，這個時間常數約為一百億年。

\begin{classroomqa}
\lecturetimestamp{00:23:46}
\audiencequestion{能否再說明一次宇宙年齡和倒數哈勃常數之間的關係？}
\lecturerresponse{冪律膨脹下可得\(H=p/t\)，所以\(H^{-1}\)與宇宙年齡相差一個約為一的因子。但在德西特空間中，\(H\)維持不變，宇宙年齡卻持續增加，因此\(H^{-1}\)代表的是指數膨脹的時間常數，而非宇宙年齡。}
\end{classroomqa}

\section{能量與宇宙視界}

如果真空能量密度在宇宙膨脹時維持不變，自然會讓人想知道真空總能量是否正在增加。這裡的回答有所限度。在廣義相對論中，能量的帳目還包含引力的部分，因此總能量的概念比初等力學中更為微妙。實際上，此處要把握的是與視界相關的要點。

在精確的德西特空間中，宇宙視界的大小固定：
\begin{equation}
D=\frac{c}{H}
=\frac{c}{\sqrt{8\pi G\rho_0/3}}.
\end{equation}
因此，視界內的體積也固定。當中的真空能量為
\begin{equation}
E_{\mathrm{hor}}=\frac{4\pi}{3}D^3\rho_0.
\end{equation}
\editorialnote{這是對00：31：06--00：32：23口頭計算的簡要整理：取視界尺度，計算球體體積，再乘以\(\rho_0\)。}
由於\(D\)和\(\rho_0\)都是常數，這個量不隨時間改變。對特定觀測者而言，在精確的德西特空間中，他所能觀測到的宇宙範圍內，真空能量維持不變。

物體也可能穿越宇宙視界，情況大致如同穿越黑洞視界。一旦到了視界之外，它們便與我們失去因果聯繫。

\begin{classroomqa}
\lecturetimestamp{00:29:27}
\audiencequestion{如果真空能量密度在宇宙膨脹時維持不變，多出來的能量從哪裡來？}
\lecturerresponse{講座以廣義相對論的觀點作答，而不是把宇宙想成牛頓力學中的盒子。在精確的德西特空間中，視界大小固定，因此永遠可被觀測到的區域內，真空能量也維持不變。}
\end{classroomqa}

\section{為何引入暴脹}

接著把同樣的德西特機制推演到遙遠的過去，只是當時的時間常數短得多，約為
\begin{equation}
H^{-1}\sim 10^{-32}\,\mathrm{s}.
\end{equation}
這就是暴脹的構想：在任何可直接觀測的時期之前，宇宙曾經歷一段指數膨脹。

基本動機在於宇宙的平坦與均勻。空間切片看來非常平坦；在大尺度上，宇宙也非常均勻。更引人注目的是，最後散射面本身就已經相當均勻，因此這種均勻化不可能只發生在後來可見的宇宙歷史中。

放癟的氣球表面皺縮不平；把它吹得夠大，任何可見的局部區域都會顯得平滑。葡萄乾或西梅也能說明相同的道理：皺縮的物體經過極大的拉伸後，在任何可見區域看來都會平滑。暴脹理論認為，宇宙正是如此變得平滑。

這也回答了此處所談的視界問題。最後散射面上相距遙遠的區域都能被我們觀測到，但在後來可見的宇宙歷史中，光沒有足夠時間從一個區域傳到另一個區域。因此，這些區域的密度和溫度幾乎一致，暗示在更早的時期，它們曾屬於同一個較小區域，之後才被大幅拉伸。

暴脹有利也有弊。它消除了宇宙對初始條件的敏感性；同樣地，也抹去了關於宇宙如何開始的許多資訊。

平坦性問題在此以更明確的形式再度出現。觀測上，曲率參數\(k\)看來為零，精確度約達百分之一或更高。在可見區域內，空間幾何看來是歐幾里得的。暴脹的解釋是，曲率半徑遠大於我們能觀測到的區域。

\begin{classroomqa}
\lecturetimestamp{00:40:05}
\audiencequestion{我們是否假設暴脹之前宇宙處於平衡狀態？}
\lecturerresponse{不是。暴脹的想法幾乎正好相反。此前不論有什麼結構，都會被大幅拉伸，使可見區域不再保留那些早期細節的資訊。}
\end{classroomqa}

\section{標量場及其位能}

為了解釋早期宇宙的有效真空能為何可能遠高於現在，可以引入一個標量場。電場和磁場都是向量，與它們不同，標量場沒有方向；它在每個位置只有一個值：
\begin{equation}
\phi=\phi(x,y,z,t).
\end{equation}
這是個假設性的場，並未對應自然界中任何已知的場，其量子也未對應任何已知粒子。

此處關鍵的能量部分是位能：
\begin{equation}
V(\phi).
\end{equation}
所謂``位能''，是指與場本身的值相關的能量，而不是與場的時間導數相關的能量。\(\phi\)取不同值時，能量也可能不同。

\begin{figure}[ht]
\centering
\includegraphics[width=0.68\textwidth]{lecture_07_figure_03.png}
\caption{定性的標量位能\(V(\phi)\)，顯示不同的場值對應不同的位能。}
\label{fig:lecture07-generic-potential}
\lectureframe{00:47:08}
\end{figure}

\begin{classroomqa}
\lecturetimestamp{00:46:34}
\audiencequestion{\(V(\phi)\)圖中的橫軸代表什麼？}
\lecturerresponse{橫軸代表場值\(\phi\)，縱軸代表位能\(V\)。}
\end{classroomqa}

如果場在空間和時間上維持定值，或在關注區域內近似定值，那麼\(V(\phi)\)的作用就與真空能量完全相同：不會稀釋。因此，先前的德西特公式只要將\(\rho_0\)替換為\(V(\phi)\)，便可沿用。特別是，
\begin{equation}
H=\sqrt{\frac{8\pi G}{3}V(\phi)}.
\end{equation}

如果\(\phi\)隨位置而變，不同區域的暴脹時間常數就會不同。但只要拉伸幅度夠大，任何可見區域都會變得近乎均勻。

\begin{classroomqa}
\lecturetimestamp{01:05:11}
\audiencequestion{這裡的\(\phi\)是否相當於時間參數？}
\lecturerresponse{在這段討論中，有限度地說是如此：場沿位能曲線演化時，它是時間的單調函數。}
\end{classroomqa}

\section{平臺、懸崖與再加熱}

這個圖景所需的位能曲線有一段非常平坦、但並非完全平坦的區域；接著到達懸崖邊緣，急遽下降，沒有完全降到零，之後又上升。

\begin{figure}[ht]
\centering
\includegraphics[width=0.68\textwidth]{lecture_07_figure_04.png}
\caption{具有高平臺、陡降區與較低極小值的暴脹位能。}
\label{fig:lecture07-inflationary-potential}
\lectureframe{01:02:30}
\end{figure}

這種形狀是高度人為設計的，但之所以提出這種模型，是因為宇宙學圖景似乎需要類似的結構。

場從平臺高處開始演化。由於斜率很小，而且膨脹所造成的有效黏滯很大，\(\phi\)變化得非常緩慢。石頭在蜂蜜中緩慢下沉，比物體在桌面上滑動更接近此處要表達的圖景。宇宙膨脹得越快，這種黏滯效應就越大。

\begin{classroomqa}
\lecturetimestamp{00:59:07}
\audiencequestion{是什麼在推動\(\phi\)演化？}
\lecturerresponse{推動它的是位能曲線的斜率。斜率很小，而宇宙學上的黏滯效應很大，所以場值變化緩慢，不會快速演化。}
\end{classroomqa}

當\(\phi\)沿著平臺緩慢演化時，\(V(\phi)\)幾乎不變。因此，宇宙的暴脹近似於真空能量維持不變時的情況。這個階段的時間常數約為
\begin{equation}
H^{-1}\sim 10^{-32}\,\mathrm{s}.
\end{equation}
口頭粗略估計，平臺高度以普朗克單位表示約為\(10^{-14}\)。\editorialnote{這是01:14:38--01:14:47口頭提出的``最佳猜測''，並非講座中推導出的數值。}

暴脹持續了多久？目前無從得知，因為平臺的寬度和坡度都不清楚。已知的是，宇宙至少膨脹了約\(e\)的五十倍，而通常引用的數字是六十倍。\editorialnote{這是對01:14:05--01:19:27口頭所談指數膨脹次數的編輯性重述。}這些次數以指數膨脹單位計：每一次都使尺度增為\(e\)倍。

當場值演化到懸崖邊緣時，位能會迅速下降。場的位能先轉化為動能，再轉化為熱和粒子，宇宙因而變得熾熱。這裡所說的``再加熱''一詞略有誤導，因為暴脹期間，先前任何熱能都已被稀釋殆盡。

較低的極小值並不在零處。該處留下的小位能，正是與當今緩慢的指數膨脹相關的能量。

\begin{classroomqa}
\lecturetimestamp{01:19:48}
\audiencequestion{平臺能否完全平坦，再由擾動以某種方式把場推離邊緣？}
\lecturerresponse{不能。位能曲線不可能完全平坦，必須帶有一定斜度。若平坦得過頭，再加熱後的宇宙就會過於不均勻。}
\end{classroomqa}

``大爆炸''在此有兩個相近的含義：可以指再加熱事件本身，也可以指最早能直接觀測到的階段，也就是最後散射面。關鍵在於，實際觀測到的一切都遠低於暴脹平臺。

\begin{classroomqa}
\lecturetimestamp{01:45:42}
\audiencequestion{再加熱本身就是大爆炸嗎？}
\lecturerresponse{你若願意，可以這麼說。但當時宇宙仍然極為熾熱且不透明。最後散射面出現得晚得多，那時宇宙已冷卻到足以變得透明。}
\end{classroomqa}

\section{量子漲落與星系的種子}

暴脹使可見宇宙變得非常均勻，但並非完全均勻。若沒有量子漲落，場會變得完全平滑；但量子力學使場在不同位置仍有微小差異。

核心圖像以縱向表示一般空間，橫向表示場值\(\phi\)。如果場完全均勻，在給定時刻，空間中每一點的\(\phi\)都會相同，而且各處會同步朝懸崖邊緣演化。如此一來，每個區域都會在同一時刻越過懸崖。量子漲落打破了這種完全同步。

兩個圖像可以說明這一點：一個是在原本均勻的場中出現小凸起；另一個是一排手挽手走向懸崖的士兵。若他們完全排齊，就會一起跌落；若有人稍微領先，有人稍微落後，他們就會在不同時刻越過懸崖。

接下來是關鍵效應。較早越過懸崖的區域會較早再加熱。等較晚的區域越過懸崖時，較早越過的區域已經膨脹了一段時間，也稍微冷卻了。這種差異會造成再加熱後的能量密度變化。

\begin{classroomqa}
\lecturetimestamp{01:34:00}
\audiencequestion{較晚再加熱的區域，比較早再加熱的區域更熱還是更冷？}
\lecturerresponse{較早再加熱的區域已有時間膨脹並冷卻。因此，較晚再加熱的區域在越過懸崖時溫度較高，這種差異也會造成密度不均勻。}
\end{classroomqa}

觀測到的效應很小：
\begin{equation}
\frac{\delta\rho}{\rho}\sim 10^{-5},
\end{equation}
約為十萬分之一。這些是原初密度漲落，並非暗能量。此處並未單憑定性圖像推導出這個數值；而是用它來限制平臺的高度與平坦程度。

這些漲落也是必要的。沒有它們，就不會有星系。再加熱之後，引力會放大初始的漲落圖樣。過密區會從欠密區吸引物質，於是過密區變得更密，欠密區則更稀，最終形成後來的星系和恆星結構。

這不只是定性的說法。可以計算膨脹宇宙中場的量子漲落，追蹤它們經過再加熱後的演化，再將結果與從星系分布重建的漲落圖樣比較。兩者符合得相當好。

這段討論也帶出熵的問題。引力造成物質結團，看來有違直覺，但並不違反熱力學第二定律。引力在統計力學中很特別：它傾向於使物質結團，而結團會提高熵，不會降低熵。

\begin{classroomqa}
\lecturetimestamp{01:40:24}
\audiencequestion{既然\(\phi\)不是電磁場，微波背景在這整個過程中如何反映這些漲落？}
\lecturerresponse{原初密度漲落的圖樣會留在最後散射面上。再加熱將場能轉化為普通粒子和光子，之後微波背景便帶有這些原初漲落留下的印記。}
\end{classroomqa}

\section{聲學振盪與下一步}

最後，討論簡短轉向聲學振盪。它們發生在後來引力造成結構成團之前。原初漲落會演化成聲學振盪，並在天空留下物理尺度已知的斑塊。

這很重要，因為在已知距離處，已知的物理尺度便能充當量尺。測量它在最後散射面上所張的角度後，就能運用幾何判斷空間是平坦、正曲率還是負曲率。詳細內容留待後續討論。

\begin{classroomqa}
\lecturetimestamp{01:47:19}
\audiencequestion{聲學振盪呢？}
\lecturerresponse{它們能在天空中充當量尺。已知的物理尺度，加上在最後散射面上觀測到的角尺度，讓我們得以測量空間幾何。詳細內容留待下一講再談。}
\end{classroomqa}